\documentclass[final]{clv2025}

\jvol{vv}
\jnum{nn}
\jyear{2025}

\dochead{Full Paper}

\pageonefooter{}

\usepackage{amsmath}
\usepackage{booktabs}
\usepackage{graphicx}
\usepackage{url}
\usepackage{microtype}
\usepackage{xcolor}

\runningtitle{Reading Without Having Read}
\runningauthor{Brew}

\begin{document}

\title{Reading Without Having Read:\\
Textual Affordances, Interpretive Frameworks,\\
and the Grounding Problem in Computational Literary Criticism}

\author{Chris Brew$^{1,2}$}

\affilblock{
    \affil{The Ohio State University$^1$\quad LexisNexis Legal \& Professional$^2$\\\quad \email{brew.2@osu.edu}}
}

\maketitle

\begin{abstract}
Can literary interpretation be decomposed into independently measurable
and manipulable components?  We present a computational framework that
generates structured multi-voice literary discussions about novels,
enabling us to isolate three factors: \emph{textual affordances} (what a
novel provides for discussion), \emph{interpretive frameworks} (the
critical perspective a reader brings), and \emph{grounding} (whether the
reader has actually read the text).  Testing across 15 Victorian and
early-modern novels (60,774 passages, 180 experimental conditions), we
find that interpretive frameworks dominate: two passage-selection
algorithms that choose almost entirely different passages (Jaccard
similarity 0.057) produce near-identical literary discussions, with
expert airtime stable to $\pm$2\% across all novels.  This is a
computational demonstration of Fish's thesis that meaning is
reader-constructed.  Without textual grounding, the system produces
\emph{pastiche}---not random hallucination but graduated borrowing from
training data, with 82\% of invented quotes drawing 90--100\% of their
vocabulary from the novel's own word stock.  The rate of faithful
quotation ranges from 57\% for canonical novels (Middlemarch) to 0\%
for obscure ones (Hester), providing an unintended instrument for
measuring an LLM's literary knowledge.  Conversational scaffolding
(host preparation) uniformly raises question density from $\sim$1 to
$\sim$11 per segment across all 15 novels, confirming that dialogue
structure is architecture-driven, not content-dependent.  We argue that
these findings connect computational text generation to long-standing
debates in reader-response theory, distant reading, and the digital
humanities.
\end{abstract}

%% ============================================================
\section{Introduction: The Computational Reader}
\label{sec:intro}

Two algorithms examine the same novel and select passages for a panel
of literary experts to discuss.  The first uses min-cost network flow
optimisation, selecting passages that satisfy each expert's thematic
demands at minimum cost.  The second uses contextual embeddings with
LLM-curated similarity, selecting passages that semantically match each
expert's profile.  The two algorithms agree on fewer than 6\% of their
selections (mean Jaccard similarity 0.057 across 15 novels).  Yet the
resulting literary discussions---multi-voice conversations in which
three experts and a host analyse the novel---are nearly
indistinguishable in character: each expert's share of airtime varies by
less than 2 percentage points regardless of which algorithm selected the
passages.

This \emph{convergence paradox} is the starting point for our
investigation.  If the texts being discussed barely overlap, but the
discussions sound the same, then the dominant factor shaping
interpretation is not the text but the reader---or, in our
computational framework, the interpretive persona.  This finding has
implications that extend well beyond podcast generation.  It provides
the first large-scale computational evidence for Stanley Fish's
influential thesis that meaning is constructed by interpretive
communities, not extracted from texts \cite{fish1980text}.  It
demonstrates that literary interpretation can be decomposed into
independently measurable components.  And it reveals what happens when
one of those components---textual grounding---is removed: the system
does not fall silent but instead produces \emph{pastiche}, a graduated
form of confabulation that recombines the novel's own vocabulary into
plausible-but-nonexistent passages.

We test nine hypotheses across 15 novels spanning 1850--1924 (four
Dickens, three Eliot, two Gaskell, one each from Collins, Forster,
Gissing, and Oliphant), totalling 60,774 enriched passages and 180
experimental conditions.  Our three main contributions are:

\begin{enumerate}
\item A demonstration that \textbf{interpretive frameworks dominate
passage selection} in determining the character of literary discussion
(Section~\ref{sec:convergence}).

\item Evidence that \textbf{ungrounded generation produces measurable
pastiche}, not random hallucination, with a confabulation rate that
serves as an unintended instrument for measuring LLM literary knowledge
(Section~\ref{sec:grounding}).

\item A practical demonstration that \textbf{literary interpretation can
be made steerable}: a scholar can explore ``what would a Marxist reading
emphasise?'' at zero marginal LLM cost
(Section~\ref{sec:editorial}).
\end{enumerate}


%% ============================================================
\section{Theoretical Framing}
\label{sec:theory}

Our framework connects to three traditions in literary theory and one in
computational creativity.

\paragraph{Interpretive communities (Fish).}
Fish~\cite{fish1980text} argued that literary texts do not contain
determinate meanings waiting to be extracted; rather, meaning is
constituted by the interpretive strategies that readers bring.  Two
readers applying different strategies to the same text will produce
different readings, not because they disagree about the text but because
they are, in effect, reading different texts.  Our system
operationalises this claim: each expert persona embodies a distinct
interpretive strategy (formalist close reading, social-historical
analysis, reader-response criticism), and the convergence paradox shows
that these strategies---not the passages selected---determine what gets
said.

\paragraph{Textual affordance (Iser).}
Iser~\cite{iser1978act} proposed that texts contain
\emph{Leerstellen}---gaps and indeterminacies---that readers must fill,
and that different texts afford different kinds of filling.  We
operationalise affordance through a seven-dimensional enrichment schema
applied to every passage in each novel: \texttt{character\_development},
\texttt{plot\_advancement}, \texttt{thematic\_depth},
\texttt{social\_critique}, \texttt{humor\_entertainment},
\texttt{atmosphere\_setting}, and \texttt{narrative\_technique}.  These
dimensions vary substantially across novels (Section~\ref{sec:affordance}),
confirming that different texts afford different interpretive
possibilities.

\paragraph{Distant reading (Moretti, Underwood).}
Moretti's~\cite{moretti2013distant} call for ``distant reading''---the
computational analysis of literary systems rather than individual
texts---has been realised primarily through topic modelling, network
analysis, and stylometry.  Our contribution is a different kind of
distant reading: we do not analyse what novels \emph{contain} but what
they \emph{afford}---what kinds of discussions they enable when filtered
through explicit critical lenses.

\paragraph{Combinational creativity (Boden).}
Boden~\cite{boden2004creative} distinguishes combinational creativity
(novel combinations of familiar ideas) from exploratory creativity
(traversal of a conceptual space) and transformational creativity
(changing the space itself).  Our system's editorial control
mechanism---in which a scholar manipulates expert demand profiles to
produce different readings---is a tool for systematic combinational
exploration of interpretive space.


%% ============================================================
\section{System Overview}
\label{sec:system}

We generate structured literary discussion podcasts through a
four-phase pipeline.  The system is the instrument, not the
contribution; we describe it briefly and refer readers to the companion
technical paper for implementation details.

\paragraph{Corpus.}  Fifteen novels (Table~\ref{tab:novels}), parsed
into passages averaging 50--80 words.  Total: 60,774 passages.

\begin{table}[t]
\centering
\small
\caption{Corpus: 15 novels, 60,774 passages.}
\label{tab:novels}
\begin{tabular}{llrr}
\toprule
Novel & Author & Year & Passages \\
\midrule
Bleak House & Dickens & 1853 & 6,916 \\
Our Mutual Friend & Dickens & 1865 & 5,775 \\
David Copperfield & Dickens & 1850 & 5,621 \\
Hard Times & Dickens & 1854 & 1,773 \\
Mill on the Floss & Eliot & 1860 & 4,654 \\
Middlemarch & Eliot & 1871 & 5,519 \\
Daniel Deronda & Eliot & 1876 & 5,170 \\
North and South & Gaskell & 1855 & 2,942 \\
Cranford & Gaskell & 1853 & 1,226 \\
A Passage to India & Forster & 1924 & 2,556 \\
No Name & Collins & 1862 & 6,057 \\
New Grub Street & Gissing & 1891 & 3,872 \\
The Odd Women & Gissing & 1893 & 3,395 \\
Miss Marjoribanks & Oliphant & 1866 & 3,019 \\
Hester & Oliphant & 1883 & 2,279 \\
\bottomrule
\end{tabular}
\end{table}

\paragraph{Phase 0: Enrichment.}  Each passage is annotated with a
20-field schema (using Claude Haiku) including an interest score
(0--5), seven \emph{provision dimensions} measuring what the passage
supplies for discussion, characters present, themes, emotional register,
and a best extractable quote.

\paragraph{Phase 1--2: Passage selection.}  We test three conditions:
\begin{itemize}
\item \textbf{Transport}: min-cost network flow matching expert demand
profiles to passage provisions.  Zero LLM tokens, $<$1 second.
\item \textbf{Embedding}: contextual embeddings with LLM-curated
similarity.  $\sim$15K Sonnet tokens, $\sim$10 seconds.
\item \textbf{No passages}: the LLM generates from prior knowledge
alone (no Phase 1--2).
\end{itemize}

\paragraph{Phase 2.5: Host preparation (optional).}  Pre-interviews
with each expert (Haiku) surface their views on assigned passages,
followed by question planning (Sonnet) that creates targeted host
questions.  This is our conversational scaffolding intervention.

\paragraph{Phase 3: Script generation.}  A single Sonnet call per
segment produces a structured podcast script with speaker turns,
quotations, and prosodic annotations.

\paragraph{Experimental design.}  180 conditions: 15 novels $\times$
2 expert panels $\times$ 3 pipelines (transport, embedding,
no-passages) $\times$ 2 host-preparation settings (with/without).
Each condition produces one complete podcast episode.


%% ============================================================
\section{The Convergence Paradox: Frameworks Dominate}
\label{sec:convergence}

\paragraph{Hypothesis DH1.}  \emph{Critical frameworks dominate passage
selection: persona explains more script variance than text.}

\paragraph{Hypothesis DH6.}  \emph{Expert prominence is
framework-stable, not text-adaptive.}

We compare the passage sets selected by transport and embedding for the
same novel and panel.  Across all 15 novels, the mean Jaccard similarity
is \textbf{0.057} (range 0.000--0.118).  Transport typically selects
49--63 passages; embedding selects 29--36.  Of these, only 3--7 are
shared.  The two methods produce nearly disjoint readings of the same
novel.

Yet the resulting scripts are similar.  Expert airtime---each expert's
share of total expert words, excluding the host---is stable across all
15 novels: Eleanor Hartley (literary critic) 32--37\%, James Blackstone
(social historian) 30--37\%, Caroline Woodcourt (close reader)
29--34\%.  The standard deviation across novels is less than 2
percentage points for each expert.

This stability is a stronger finding than we predicted.  Our original
hypothesis (H6 in the companion paper) predicted that expert prominence
would \emph{adapt} to textual affordance---that a novel rich in social
critique would give the social historian more airtime.  The data shows
the opposite: the frameworks are \emph{robust}, maintaining their
character regardless of what the novel provides.  For digital humanities,
this is the more interesting result.  It means that the interpretive
frameworks function as genuine critical lenses---consistent heuristics
for reading, not chameleons that reshape themselves to fit each text.

\subsection{Two Reading Strategies, One Discussion}

The two selection methods do not merely pick different passages; they
embody recognisably different \emph{reading strategies}.  Transport
selects nearly twice as many passages (59 vs.\ 32 per run), drawn from
40\% more chapters (28 vs.\ 20), with higher mean interest scores (4.44
vs.\ 4.01).  Examining the provision-dimension profiles of selected
passages reveals a systematic difference: transport favours
\texttt{character\_development} (+8.3~pp) and
\texttt{plot\_advancement} (+10.2~pp), while embedding favours
\texttt{social\_critique} (+13.4~pp) and \texttt{thematic\_depth}
(+6.9~pp).

In literary-pedagogical terms, transport produces a \emph{synoptic}
reading---broad coverage of the novel's story and characters---while
embedding produces an \emph{analytical} reading---concentrated
engagement with the novel's arguments and themes.  These are familiar
strategies in the seminar room, and the fact that they emerge from
purely algorithmic choices (min-cost flow vs.\ semantic similarity) is
itself noteworthy.

Yet the resulting discussions converge.  This is the force of the
convergence paradox: not merely that different passages lead to similar
discussions, but that \emph{different reading strategies}---synoptic
vs.\ analytical, breadth vs.\ depth---produce discussions whose
character is determined by the interpretive framework, not by the
reading strategy.

\paragraph{Interpretation.}  The convergence paradox operationalises
Fish's thesis at scale.  When two recognisably different reading
strategies produce convergent discussions, the dominant variable is the
interpretive community (the expert panel), not the mode of reading.
A Marxist close reading and a Marxist survey reading of the same novel
differ in their evidence base but not in their analytical character.
The reader makes the reading.


%% ============================================================
\section{Textual Affordances as Novel Fingerprints}
\label{sec:affordance}

\paragraph{Hypothesis DH2.}  \emph{Textual affordances are measurable
and vary: provision dimensions discriminate novels.}

\paragraph{Hypothesis DH7.}  \emph{The enrichment schema partially
transfers across novels.}

Our seven provision dimensions produce distinctive profiles for each
novel.  The most diagnostic dimension is
\texttt{humor\_entertainment}, which has a coefficient of variation of
0.704 across the 15 novels: \textit{Miss Marjoribanks} (Oliphant) has
the richest comedic supply, while \textit{New Grub Street} (Gissing)
and \textit{The Odd Women} (Gissing) are among the sparsest.

Six of the seven dimensions transfer across all 15 novels without
modification.  The exception is \texttt{atmosphere\_setting}, which is
uniformly low ($<$6\% of passages rated strong) across all novels.  We
interpret this as a bias in the enrichment schema toward features that
are diagnostic in Realist fiction---the schema was designed on
\textit{Bleak House} and reflects Dickensian priorities.  The dimension
is not absent from these novels; it is present but the schema does not
capture it well.

\paragraph{Affordance profiles.}  Each novel's provision-dimension
profile constitutes a measurable \emph{affordance fingerprint}.
\textit{Hard Times} affords social critique (20.2\% strong) more than
any other novel.  \textit{Miss Marjoribanks} affords comedy (27.1\%)
and character study (40.5\%).  These are not subjective judgements but
reproducible computational measurements, and they constrain (without
determining) the discussions that can be generated.


%% ============================================================
\section{Reading Without Having Read}
\label{sec:grounding}

\paragraph{Hypothesis DH3.}  \emph{Grounding is prerequisite for
faithful interpretation: the gap between grounded and ungrounded
generation is $\geq$40 percentage points.}

\paragraph{Hypothesis DH4.}  \emph{Confabulation measures training
exposure: ungrounded verification rates correlate with novel prominence.}

\paragraph{Hypothesis DH5.}  \emph{Confabulation is pastiche, not
random hallucination.}

The grounding gap is the central empirical finding of this study.  When
given source passages (transport or embedding), the system quotes the
novel accurately: mean verification rates of 98.0\% (transport) and
97.2\% (embedding).  When generating from prior knowledge alone
(no-passages condition), the mean drops to 48.9\%, with a range of
0--57\% depending on the novel.

\subsection{The Knowledge Gradient}

Table~\ref{tab:nop_verification} shows the ungrounded verification
rate for each novel, ordered by decreasing accuracy.  The pattern is
clear: the LLM quotes canonical, widely-read novels more accurately
than obscure ones.

\begin{table}[t]
\centering
\small
\caption{Ungrounded (no-passages) quote verification by novel.
Verification uses 5-tier matching against the full enriched corpus:
\emph{verified} ($\geq$60\% match), \emph{paraphrase} (30--59\%),
\emph{distant echo} (10--29\%), \emph{no clear source} (1--9\%),
\emph{invented} (0\%).}
\label{tab:nop_verification}
\begin{tabular}{llrr}
\toprule
Novel & Author & Verified\% & Invented\% \\
\midrule
Middlemarch & Eliot & 57 & 21 \\
Bleak House & Dickens & 53 & 22 \\
David Copperfield & Dickens & 44 & 17 \\
Hard Times & Dickens & 39 & 40 \\
Passage to India & Forster & 39 & 34 \\
Mill on the Floss & Eliot & 37 & 32 \\
Cranford & Gaskell & 33 & 61 \\
Daniel Deronda & Eliot & 30 & 36 \\
Our Mutual Friend & Dickens & 29 & 55 \\
New Grub Street & Gissing & 20 & 65 \\
Odd Women & Gissing & 6 & 91 \\
Miss Marjoribanks & Oliphant & 3 & 44 \\
North and South & Gaskell & 2 & 77 \\
Hester & Oliphant & 0 & 78 \\
No Name & Collins & 0 & 100 \\
\bottomrule
\end{tabular}
\end{table}

The grounding gap ranges from approximately 30 percentage points
(Middlemarch: grounded 99\%, ungrounded 57\%) to nearly 100 percentage
points (Hester, No Name: grounded $\sim$98\%, ungrounded 0\%).
\textbf{Hypothesis DH3 is confirmed with a much larger effect than
predicted}, and \textbf{DH4 is confirmed}: novel prominence is a strong
predictor of ungrounded accuracy.

\subsection{Pastiche, Not Hallucination}

The most distinctive finding concerns the \emph{nature} of
confabulation.  We performed a vocabulary autopsy on all 726 invented
quotes (0\% match with any passage in the novel).  The results are
striking:

\begin{itemize}
\item \textbf{82\% of invented quotes} draw 90--100\% of their
individual words from the novel's own vocabulary.
\item The mean percentage of quote words appearing somewhere in the
novel is \textbf{95.1\%}.
\item The minimum is 67\%---even the worst confabulations are
constructed from the novel's word stock.
\end{itemize}

This is not random generation.  The LLM is writing \emph{in the style
of} the novel rather than \emph{from} the novel.  It recombines real
vocabulary into plausible-but-nonexistent sequences.  In literary-critical
terms, this is \emph{pastiche}---imitation of an author's manner
without direct quotation from their text.

For the DH community, this reframes the ``hallucination problem'' in
LLM-generated literary discussion.  When the system lacks grounding, it
does not produce gibberish; it produces a kind of creative
imitation---sometimes quite convincing---that draws on whatever the
model absorbed during training.  The graduated nature of this
imitation (from perfect quotation of famous passages to pure pastiche
of obscure ones) makes it a measurable phenomenon, not merely an error
to be eliminated.

\paragraph{Interpretation.}  We suggest that confabulation rate
functions as an unintended \textbf{instrument for measuring LLM literary
knowledge}.  A novel that the model can quote accurately from memory
(Middlemarch, 57\% verified) is one the model has deeply absorbed.  A
novel it cannot quote at all (Hester, 0\%) is one it has encountered
only glancingly, or not at all.  This is, to our knowledge, the first
proposal to use confabulation rate as a measurement tool rather than
treating it solely as a failure mode.


%% ============================================================
\section{Conversational Scaffolding}
\label{sec:scaffolding}

\paragraph{Hypothesis DH8.}  \emph{Conversational structure is
scaffolding-driven, not content-driven.}

Host preparation---a phase in which the system conducts pre-interviews
with each expert persona and then plans targeted questions---transforms
the generated discussions.  Without host preparation, the system produces
extended expert monologues with a mean of $\sim$1 question per segment.
With host preparation, the mean rises to $\sim$11 questions per segment.

Crucially, this effect is \textbf{uniform across all 15 novels}.
The ratio of questions with host prep to questions without ranges from
5.3$\times$ (\textit{Our Mutual Friend}) to 14.1$\times$
(\textit{Cranford}), but every novel shows the transformation from
monologue to dialogue.  This uniformity confirms that the effect is
architectural---it is a property of the conversational scaffolding, not
of the literary content being discussed.

\paragraph{Interpretation.}  For digital humanities, this finding
separates two concerns that are often conflated: the \emph{substance}
of a literary discussion (what is said about the novel) and its
\emph{dynamics} (how speakers interact).  The dynamics can be designed
independently of the substance, through structural interventions that
operate at the level of conversation architecture rather than literary
content.


%% ============================================================
\section{Editorial Control}
\label{sec:editorial}

\paragraph{Hypothesis DH9.}  \emph{Editorial manipulation is tractable:
demand-profile changes produce predictable, graduated passage shifts.}

The transport selection method formulates passage selection as a
min-cost flow problem in which each expert has a \emph{demand profile}
over the seven provision dimensions.  By modifying these demand
profiles---increasing one expert's appetite for social critique,
another's for narrative technique---a scholar can explore how different
critical emphases reshape the discussion.

We tested this by varying the demand profiles across multiple
configurations, from conservative adjustments (single-dimension shifts)
to radical panel replacement (all three experts swapped).  The resulting
passage-set Jaccard similarities show a \textbf{graduated response}:
conservative shifts produce Jaccard values of 0.7--0.8 with the
baseline, while full panel replacement produces 0.5--0.6.  The
editorial intention is legible as a predictable flow of passage changes.

This exploration incurs \textbf{zero marginal LLM cost}.  The min-cost
flow solver runs in $<$100ms with no API calls.  A scholar can preview
dozens of configurations before committing to the expensive script
generation step.  This ``preview-before-commit'' workflow is unique
to the transport formulation; embedding-based selection offers no
equivalent mechanism for systematic exploration.

\paragraph{Interpretation.}  In Boden's~\cite{boden2004creative}
terms, this is a tool for \emph{exploratory creativity}---systematic
traversal of an interpretive space defined by provision dimensions and
expert demands.  The scholar does not create new critical lenses but
explores the consequences of adjusting existing ones, at a granularity
and speed that would be impractical with human panels.


%% ============================================================
\section{Discussion}
\label{sec:discussion}

\subsection{What Does It Mean to ``Operationalise'' Interpretation?}

Our framework does not claim to replicate human literary interpretation.
It claims something more modest and, we believe, more useful: that the
\emph{structure} of interpretation---the relationship between text,
reader, and resulting discourse---can be modelled as a set of
independently variable factors, each of which can be measured,
manipulated, and compared.

The convergence paradox (Section~\ref{sec:convergence}) suggests that
the most powerful factor is the interpretive framework.  The grounding
results (Section~\ref{sec:grounding}) show that access to the actual
text is a prerequisite for accuracy but not for fluency.  The
scaffolding results (Section~\ref{sec:scaffolding}) demonstrate that
conversational dynamics can be designed independently of content.  These
are not literary-critical insights in themselves; they are
\emph{structural} insights about how literary-critical discourse is
produced.

\subsection{The Schema Problem}

Our enrichment schema---seven provision dimensions derived from
\textit{Bleak House}---partially transfers to 14 additional novels.
The failure of \texttt{atmosphere\_setting} suggests that the schema is
inflected by the tradition in which it was designed.  A schema developed
from modernist fiction, or from non-English literature, would likely
capture different affordances.  We regard this as a productive
limitation: the fact that schemas are perspective-dependent is itself a
finding about the nature of literary description.

\subsection{Limitations}

\textbf{Automated metrics only.}  We report verification rates,
airtime distributions, and structural statistics, but no human
evaluation of discussion quality, insight, or literary value.

\textbf{English canon.}  All 15 novels are in English, from
1850--1924.  We cannot claim generalisability to other languages,
periods, or genres.

\textbf{Enrichment bias.}  The provision dimensions reflect the
priorities of Victorian Realism.  Application to other traditions would
require schema adaptation.

\textbf{No causal claims.}  Our experimental design supports
correlational findings (e.g., novel prominence correlates with
verification rate) but does not establish causal mechanisms.


%% ============================================================
\section{Related Work}
\label{sec:related}

\paragraph{Computational literary criticism.}
Moretti's~\cite{moretti2013distant} ``distant reading'' programme and
Underwood's~\cite{underwood2019distant} statistical analyses of
literary history have established computational methods in literary
studies.  Our work differs in generating \emph{discourse about}
literature rather than analysing literature directly.

\paragraph{AI-generated conversation.}
Google's NotebookLM~\cite{notebooklm2024} generates audio discussions
from uploaded documents.  Our contribution is the explicit
decomposition of the generation process into manipulable components,
enabling controlled comparison across conditions.

\paragraph{Retrieval-augmented generation.}
Lewis et al.~\cite{lewis2020rag} and Gao et al.~\cite{gao2024rag}
establish RAG as the standard approach to grounding LLM output.  Our
finding that grounding method matters far less than the fact of
grounding extends this literature.

\paragraph{Hallucination and confabulation.}
Our pastiche finding---that confabulation is structured, graduated, and
vocabulary-preserving---extends work on LLM hallucination by
characterising confabulation as a literary phenomenon (pastiche) rather
than solely a technical failure.


%% ============================================================
\section{Conclusion}
\label{sec:conclusion}

We have shown that literary interpretation, modelled as a computational
process, decomposes into three independently measurable factors:
interpretive frameworks, textual affordances, and grounding.  Of these,
frameworks dominate---a finding that provides the first large-scale
computational evidence for Fish's thesis that meaning is
reader-constructed.

When grounding is removed, the system does not fail randomly but
produces graduated pastiche, constructing plausible quotations from
the novel's own vocabulary.  This pastiche is more faithful for
canonical novels and degrades to pure invention for obscure ones,
providing a novel instrument for measuring LLM literary knowledge.

Conversational scaffolding transforms monologue into dialogue
uniformly across all 15 novels, confirming that discussion dynamics
can be designed independently of literary content.  Editorial control
through demand-profile manipulation enables zero-cost exploration of
interpretive space.

These findings connect AI text generation to long-standing debates in
reader-response theory and suggest that computational frameworks can
contribute to---not replace---the humanities' understanding of how
reading works.


\bibliography{refs}
\bibliographystyle{compling}

\end{document}
